% Einheit 1 von 4 – Negative Zahlen kennen und ordnen (bank/rationale-zahlen/e1.jsonl, zusammenbau v0.1)
%% TODO Zweigzeile Teil 1: was hier gelernt wird (Satz in Schülersprache aus dem Einheitstitel) – Einheit 1
\einheitenkopf[e1]{Einheit 1 von 4 · Negative Zahlen kennen und ordnen}
\zweigzeile{ab Kl. 6, je nach Buch bis Kl. 7 · P10}
\setcounter{aufgabe}{7}

%% TODO Titel: Ich-kann-Satz fehlt in der Bank (Platzhalter: Kettenname) – Nr. 8
%% TODO Anweisung: ein Satz über der ersten Teilaufgabe fehlt in der Bank – Nr. 8
\begin{aufgabe}{Ordnen}
%% TODO \rechenplatz{n}: Schrittzahl des Grundfalls fehlt in der Bank (nur bei fester Schreibform, 2.3 b)
\begin{teile}
\teil Schreibe an jeden Strich seine Zahl. Abstand zwischen zwei Strichen: $1$.

\zahlenstrahl[xmin=-5,xmax=5,xstep=1]{0/0}
\teil Trage ein: $-4$, $3$ und $-8$.

\zahlenstrahl[xmin=-10,xmax=10,xstep=1]{}
\teil $-17$ – Gegenzahl und Betrag? Gegenzahl \leerfeld , Betrag \leerfeld
\teil $-15$ oder $-9$ – welche ist kleiner? \leerfeld
\teil Trage ein: $-2{,}5$, $0{,}5$ und $-\frac{3}{2}$.

\zahlenstrahl[xmin=-3,xmax=3,xstep=0.5,karo=1]{}
\teil Ordne aufsteigend: $-0{,}8$; $\frac{1}{4}$; $-1{,}2$; $-\frac{3}{4}$ \leerfeld $<$ \leerfeld $<$ \leerfeld $<$ \leerfeld
\teil Gib eine Zahl an, die größer ist als $-320$. \hfill (P10 2015 OS) \\ \leerfeld
\teil Ordne aufsteigend: $-\frac{2}{5}$; $1{,}3$; $-0{,}43$; $0{,}9$ \hfill (P10 2014 OS) \\ \leerfeld $<$ \leerfeld $<$ \leerfeld $<$ \leerfeld
\end{teile}
\end{aufgabe}

%% TODO Titel: Ich-kann-Satz fehlt in der Bank (Platzhalter: Kettenname) – Nr. 9
%% TODO Anweisung: ein Satz über der ersten Teilaufgabe fehlt in der Bank – Nr. 9
\begin{aufgabe}{Mitte zweier Zahlen}
\begin{teile}
\teil Mitte von $-8$ und $2$? \leerfeld
\end{teile}
\end{aufgabe}

%% TODO Titel: Ich-kann-Satz fehlt in der Bank (Platzhalter: Kettenname) – Nr. 10
%% TODO Anweisung: ein Satz über der ersten Teilaufgabe fehlt in der Bank – Nr. 10
\begin{aufgabe}{runden}
\begin{teile}
\teil $-3{,}47$ auf Zehntel runden \leerfeld
\end{teile}
\end{aufgabe}

%% TODO Titel: Ich-kann-Satz fehlt in der Bank (Platzhalter: Kettenname) – Nr. 11
%% TODO Anweisung: ein Satz über der ersten Teilaufgabe fehlt in der Bank – Nr. 11
\begin{aufgabe}{Punkte in vier Quadranten}
\begin{teile}
\teil Punkt P – welche Koordinaten? Lies ab.

\begin{ksys}[xmin=-5,xmax=5,ymin=-5,ymax=5,ablesen] \punkt{-3}{2}{P} \end{ksys}
P( \leerfeld | \leerfeld )
\end{teile}
\end{aufgabe}

%% TODO Titel: Ich-kann-Satz fehlt in der Bank (Platzhalter: Kettenname) – Nr. 12
%% TODO Anweisung: ein Satz über der ersten Teilaufgabe fehlt in der Bank – Nr. 12
\begin{aufgabe}{Ordnen – Fehler finden, Begründen, Darstellungswechsel, Anwendung}
\begin{teile}
\teil Tim schreibt: $-45 > -30$, weil $45 > 30$. Finde den Fehler.
\teil Warum ist $-17$ kleiner als $-4$? Begründe.
\teil Ein Taucher ist $18$ m unter dem Meeresspiegel. Welche Zahl? \leerfeld[m]
\teil Tiefste Punkte an Land: Kaspisches Meer $-28$ m, Totes Meer $-430$ m, Death Valley $-86$ m, Qattara-Senke $-133$ m. Welcher Ort liegt am tiefsten?
\end{teile}
\end{aufgabe}
